\begin{exercise}{3.4.4}
Call {\sffamily LC} Configuration $\langle P,s \rangle$ \textit{sensible} if the set of locations on which $s$ is defines, $dom(s)$, contains all the locations that occur in the phrase $P$. Prove by induction on the structure of $P$ that if $\langle P,s \rangle$ is \textit{sensible}, than it is not stuck. Prove by Rule Induction for $\rightarrow$ that if $\langle P,s \rangle \rightarrow \langle P',s' \rangle$ and $\langle P,s \rangle$ is \textit{sensible}, then so is $\langle P',s' \rangle$ and $dom(s') = dom(s)$. Deduce that a stuck configuration can never be reached by a series of transitions from a sensible configuration.

\end{exercise}

\begin{solution}
\begin{enumerate}
    \item \textbf{Definitions}. We define:
    \begin{itemize}
        \item $\langle P,s \rangle$ \textit{sensible} if the set of locations on which $s$ is defines, $dom(s)$, contains all the locations that occur in the phrase $P$;
        \item $\langle P,s \rangle$ \textit{stuck} if it is not terminal and there is no possible transition $\langle P,s \rangle \rightarrow \langle P',s' \rangle$.
    \end{itemize}
    \item \textbf{By structural induction, \textit{sensible} $\implies$ not \textit{stuck}}. \\
    We want to show that if $\langle P,s \rangle$ is \textit{sensible}, then it is not \textit{stuck}. We proceed by structural induction on $P$.
    \begin{enumerate}
        \item \textsc{Base cases} 
        \begin{enumerate}
            \item \textbf{Terminal configurations} ($P=n$ (integer), $P=b$ (boolean value) or $P=skip$). By definition $P$ is terminal and a \textit{stuck} configuration is by definition not terminal, thus $P$ is not \textit{stuck};
            \item \textbf{Memory lookup} ($P=!\ell $). If $\langle !\ell, s\rangle$ is sensible then $\ell \in dom(s)$ hence by 
            \[(\overset{\rightarrow}{loc} ) \quad \frac{}{\langle !\ell ,s \rangle \to \langle n,s\rangle} \text{ with }  s(\ell) = n\]
            there is a transition. Since a such a configuration exists, $\langle !\ell, s\rangle$ is not stuck.
        \end{enumerate}
        \item \textsc{Inductive step}. Assume by inductive hypothesis (IH) that the property we're studying holds for $P_1$ and $P_2$.
            \begin{enumerate}
                \item \textbf{Binary operations} ($P=P_1$ $op$ $P_2$)
                \begin{enumerate}
                    \item if $P_1$ is not a value, then by IH there is a configuration $\langle P_1,s\rangle \to \langle P_1', s'\rangle$. By the rule 
                    \[ (\overset{\rightarrow}{op_1}) \quad \frac{\langle E_1,s\rangle \to \langle E_1', s'\rangle}{\langle E_1 \, op \, E_2,s\rangle \to \langle E_1' \,op\, E_2, s'\rangle}\]
                    there is a configuration for $P$, so it is not \textit{stuck};
                    \item if $P_1$ is a value, but $P_2$ isn't the by IH there is a configuration $\langle P_2,s\rangle \to \langle P_2', s'\rangle$. By the rule 
                    \[ (\overset{\rightarrow}{op_2}) \quad \frac{\langle E_2,s\rangle \to \langle E_2', s'\rangle}{\langle n \, op \, E_2,s\rangle \to \langle n \,op\, E_2', s'\rangle}\]
                    there is a configuration for $P$, so it is not \textit{stuck};
                    \item if both $P_1$ and $P_2$ are values, by the rule 
                    \[ (\overset{\rightarrow}{op_3}) \quad \frac{}{\langle n \, op \, m,s\rangle \to \langle c, s\rangle} \; \text{with } c = n \, op \, m\]
                    there is a configuration for $P$, so it is not \textit{stuck};
                \end{enumerate}
                \item \textbf{Assignment} ($P=\ell :=P_1)$
                \begin{enumerate}
                    \item if $P_1$ is not a value, then by IH there is a configuration $\langle P_1,s\rangle \to \langle P_1', s'\rangle$. By the rule 
                    \[ (\overset{\rightarrow}{set_1}) \quad \frac{\langle E,s\rangle \to \langle E', s'\rangle}{\langle \ell:=E,s\rangle \to \langle \ell:=E', s'\rangle}\]
                    there is a configuration for $P$, so it is not \textit{stuck};
                    \item if $P_1$ is a value $n$ since $\ell \in dom(s)$ we can use the rule 
                    \[ (\overset{\rightarrow}{set_2}) \quad \frac{}{\langle \ell:=n,s\rangle \to \langle skip, s[\ell \mapsto n]\rangle}\]
                    without the risk of a \textit{stuck} configuration, then there is a configuration for $P$, so it is not \textit{stuck};
                \end{enumerate}
                \item \textbf{Composition} ($P=P_1 ; P_2)$
                \begin{enumerate}
                    \item if $P_1$ is not a skip, then by IH there is a configuration $\langle P_1,s\rangle \to \langle P_1', s'\rangle$. By the rule 
                    \[ (\overset{\rightarrow}{seq_1}) \quad \frac{\langle C_1,s\rangle \to \langle C_1', s'\rangle}{\langle C_1;C_2,s\rangle \to \langle C_1';C_2, s'\rangle}\]
                    there is a configuration for $P$, so it is not \textit{stuck};
                    \item if $P_1$ is a value, then by the rule 
                    \[ (\overset{\rightarrow}{seq_2}) \quad \frac{}{\langle skip;C_2,s\rangle \to \langle C_2,s\rangle}\]
                    there is a configuration for $P$, so it is not \textit{stuck};
                \end{enumerate}
            \item \textbf{If-Then-Else} ($P=\text{if }B \text{ then } C_1 \text{ else }C_2)$
                We assume $\langle P,s\rangle = \langle \text{if }B \text{ then } C_1 \text{ else }C_2, s \rangle$ to be sensible, i.e. $Loc(B) \cup Loc(C_1) \cup Loc(C_2) \subseteq dom(s)$, so in particular 
                \begin{itemize}
                    \item $\langle B,s\rangle$ is sensible; 
                    \item $\langle C_1,s\rangle$ is sensible; 
                    \item $\langle C_2,s\rangle$ is sensible.
                \end{itemize}
                We can have three cases on the boolean expression $B$:
                \begin{enumerate}
                    \item $B$ is not yet a value, then we can apply 
                      \[ (\overset{\rightarrow}{if1}) \quad \frac{\langle B,s\rangle \to \langle B', s'\rangle}{\langle \text{if }B \text{ then }C_1 \text{ else }C_2, s\rangle \to \langle \text{if }B' \text{ then }C_1 \text{ else }C_2, s'\rangle}\]
                    \item $B = true$, then we can apply 
                     \[ (\overset{\rightarrow}{if2}) \quad \frac{}{\langle \text{if }true \text{ then }C_1 \text{ else }C_2, s\rangle \to \langle C_1,s\rangle}\]
                    \item $B = false$, then we can apply \[ (\overset{\rightarrow}{if3}) \quad \frac{}{\langle \text{if }false \text{ then }C_1 \text{ else }C_2, s\rangle \to \langle C_2,s\rangle}\]
                \end{enumerate}
                Hence in all cases a transition exists, therefore $\langle P,s\rangle $ is not stuck.


                   \item \textbf{While} ($P=\text{while }B \text{ do } C$
                We assume $\langle P,s\rangle = \langle \text{while }B \text{ do } C, s \rangle$ to be sensible, i.e. $Loc(B) \cup Loc(C)  \subseteq dom(s)$.
                A while loop is never evaluated directly, instead it unfolds into an if statement: 
                \[ (\overset{\rightarrow}{whl}) \quad \frac{}{\langle \text{while }B \text{ do } C, s\rangle \to \langle \text{if }B \text{ then }(C;\text{while }B \text{ do } C) \text{ else }skip, s\rangle}\]
                So a transition always exists, thus  $\langle P,s\rangle $ is not stuck.
                \end{enumerate}

    \end{enumerate}
    \item \textbf{By rule induction, preservation of \textit{sensible} configuration}. \\
    We want to prove that if $\langle P,s \rangle \rightarrow \langle P',s' \rangle$ and $\langle P,s \rangle$ is \textit{sensible}, then so is $\langle P',s' \rangle$ and $dom(s') = dom(s)$.
    \begin{enumerate}
        \item \textsc{Closure under axioms}
        \begin{enumerate}
            \item \textbf{Memory lookup}
            \[(\overset{\rightarrow}{loc} ) \quad \frac{}{\langle !\ell ,s \rangle \to \langle n,s\rangle} \text{ with }  s(\ell) = n\]
            Note that the memory state doesn't change $s'=s$, so $dom(s') = dom(s)$. Then $s(\ell)$ is an integer which doesn't contain locations ($Loc(s(\ell))= \emptyset$) and the empty set is a subset of every domain, hence the resulting configuration is \textit{sensible}.
            \item \textbf{Assignment}
             \[ (\overset{\rightarrow}{set_2}) \quad \frac{}{\langle \ell:=n,s\rangle \to \langle skip, s[\ell \mapsto n]\rangle}\]
             Since the starting configuration is \textit{sensible} $\ell \in dom(s)$, hence $s[\ell \mapsto n]$ changes the value of an existing location without changing the set of locations, so $dom(s') = dom(s)$.
            \item \textbf{Integer operation}
             \[ (\overset{\rightarrow}{op_3}) \quad \frac{}{\langle n \, op \, m,s\rangle \to \langle c, s\rangle} \; \text{with } c = n \, op \, m\]
              Note that the memory state doesn't change $s'=s$, so $dom(s') = dom(s)$. Then $c$ is an integer which doesn't contain locations ($Loc(c)= \emptyset$) and the empty set is a subset of every domain, hence the resulting configuration is \textit{sensible}.
              \item \textbf{Skip}
             \[ (\overset{\rightarrow}{seq_2}) \quad \frac{}{\langle skip;C_2,s\rangle \to \langle C_2,s\rangle}\]
              By hypothesis $\langle skip;C_2,s\rangle$ is sensible, i.e. $Loc(skip)=\emptyset \cup Loc(C_2) \subseteq dom(s) $. Since the memory state doesn't change $s'=s$, so $dom(s') = dom(s)$, hence $Loc(C_2) \subseteq dom(s') $ so the resulting configuration is \textit{sensible}.
               \item \textbf{If-True}
               \[ (\overset{\rightarrow}{if2}) \quad \frac{}{\langle \text{if }true \text{ then }C_1 \text{ else }C_2, s\rangle \to \langle C_1,s\rangle}\]
              By hypothesis $\langle \text{if }true \text{ then }C_1 \text{ else }C_2, s\rangle$ is sensible, i.e. $Loc(B) \cup Loc(C_1) \cup Loc(C_2) \subseteq dom(s)$. Since the memory state doesn't change $s'=s$, so $dom(s') = dom(s)$, hence $Loc(B) \cup Loc(C_1) \cup Loc(C_2) \subseteq dom(s') $ so the resulting configuration is \textit{sensible}.
               \item \textbf{If-False}
              \[ (\overset{\rightarrow}{if3}) \quad \frac{}{\langle \text{if }false \text{ then }C_1 \text{ else }C_2, s\rangle \to \langle C_2,s\rangle}\]
              By hypothesis $\langle \text{if }false \text{ then }C_1 \text{ else }C_2, s\rangle$ is sensible, i.e. $Loc(B) \cup Loc(C_1) \cup Loc(C_2) \subseteq dom(s)$. Since the memory state doesn't change $s'=s$, so $dom(s') = dom(s)$, hence $Loc(B) \cup Loc(C_1) \cup Loc(C_2) \subseteq dom(s') $ so the resulting configuration is \textit{sensible}.
               \item \textbf{While}
             \[ (\overset{\rightarrow}{whl}) \quad \frac{}{\langle \text{while }B \text{ do } C, s\rangle \to \langle \text{if }B \text{ then }(\text{while }B \text{ do } C)) \text{ else }skip, s\rangle}\]
              By hypothesis $\text{while }B \text{ do } C, s\rangle$ is sensible, i.e. $Loc(B) \cup Loc(C) \subseteq dom(s)$. Since the memory state doesn't change $s'=s$, so $dom(s') = dom(s)$, hence $Loc(B) \cup Loc(C) \subseteq dom(s') $ so the resulting configuration is \textit{sensible}.
        \end{enumerate}
        \item \textsc{Closure under rules} 
          \begin{enumerate}
                \item \textbf{Binary operation}
                \[ (\overset{\rightarrow}{op_1}) \quad \frac{\langle E_1,s\rangle \to \langle E_1', s'\rangle}{\langle E_1 \, op \, E_2,s\rangle \to \langle E_1' \,op\, E_2, s'\rangle}\]
                Assume $\langle E_1 \, op \, E_2,s\rangle $ \textit{sensible}, i.e. $Loc(E_1) \cup Loc(E_2) \subseteq dom(s)$.
                \\
                By Inductive Hypothesis (IH) for $\langle E_1,s\rangle \to \langle E_1', s'\rangle$ we have that $\langle E_1', s'\rangle$ is \textit{sensible}, i.e. $Loc(E_1') \subseteq dom(s') $ and $dom(s') = dom(s)$.
                \\
                Since $Loc(E_2) \subseteq dom(s) $ and $dom(s') = dom(s)$, $Loc(E_2) \subseteq dom(s') $.
                \\
                Hence $Loc(E_2) \subseteq dom(s') $ and $Loc(E_1') \subseteq dom(s') $ thus $Loc(E_1') \cup Loc(E_2) \subseteq dom(s') $ so the resulting configuration is \textit{sensible}.
                \item \textbf{Assignment} 
                 \[ (\overset{\rightarrow}{set_1}) \quad \frac{\langle E,s\rangle \to \langle E', s'\rangle}{\langle \ell:=E,s\rangle \to \langle \ell:=E', s'\rangle}\]
                Assume $\langle \ell:=E,s\rangle $ is \textit{sensible}, i.e. $Loc(\ell:=E) = \{\ell\}\cup Loc(E) \subseteq dom(s)$.
                By Inductive Hypothesis (IH) for $\langle E,s\rangle \to \langle E', s'\rangle$ we have that $\langle E', s'\rangle$ is \textit{sensible}, i.e. $Loc(E') \subseteq dom(s') $ and $dom(s') = dom(s)$.
                \\
                Since $Loc(\ell) \subseteq dom(s) $ and $dom(s') = dom(s)$, $Loc(\ell) \subseteq dom(s') $.
                \\
                Hence $Loc(\ell) \subseteq dom(s') $ and $Loc(E') \subseteq dom(s') $ thus $Loc(\ell) \cup Loc(E') \subseteq dom(s') $ so the resulting configuration is \textit{sensible}.
                \item \textbf{Sequence} 
                 \[ (\overset{\rightarrow}{seq_1}) \quad \frac{\langle C_1,s\rangle \to \langle C_1', s'\rangle}{\langle C_1;C_2,s\rangle \to \langle C_1';C_2, s'\rangle}\]
                 Assume $\langle C_1;C_2,s\rangle $ \textit{sensible}, i.e. $Loc(C_1) \cup Loc(C_2) \subseteq dom(s)$.
                \\
                By Inductive Hypothesis (IH) for $\langle C_1,s\rangle \to \langle C_1', s'\rangle$ we have that $\langle C_1', s'\rangle$ is \textit{sensible}, i.e. $Loc(C_1') \subseteq dom(s') $ and $dom(s') = dom(s)$.
                \\
                Since $Loc(C_2) \subseteq dom(s) $ and $dom(s') = dom(s)$, $Loc(C_2) \subseteq dom(s') $.
                \\
                Hence $Loc(C_2) \subseteq dom(s') $ and $Loc(C_1') \subseteq dom(s') $ thus $Loc(C_1') \cup Loc(C_2) \subseteq dom(s') $ so the resulting configuration is \textit{sensible}.
                \item \textbf{If-Then-else} 
                 \[ (\overset{\rightarrow}{if1}) \quad \frac{\langle B,s\rangle \to \langle B', s'\rangle}{\langle \text{if }B \text{ then }C_1 \text{ else }C_2, s\rangle \to \langle \text{if }B' \text{ then }C_1 \text{ else }C_2, s'\rangle}\]
                 Assume $\langle \text{if }B \text{ then }C_1 \text{ else }C_2, s\rangle$ \textit{sensible}, i.e. $Loc(B) \cup Loc(C_1) \cup Loc(C_2) \subseteq dom(s)$.
                \\
                By Inductive Hypothesis (IH) for $\langle B,s\rangle \to \langle B', s'\rangle$ we have that $\langle B', s'\rangle$ is \textit{sensible}, i.e. $Loc(B') \subseteq dom(s') $ and $dom(s') = dom(s)$.
                \\
                Since $Loc(C_1) \cup Loc(C_2) \subseteq dom(s) $ and $dom(s') = dom(s)$, $Loc(C_1) \cup Loc(C_2)\subseteq dom(s') $.
                \\
                Hence $Loc(C_1) \cup Loc(C_2)\subseteq dom(s')$ and $Loc(B') \subseteq dom(s') $ thus $Loc(C_1') \cup Loc(C_2) \cup Loc(B') \subseteq dom(s') = dom(s)$ so the resulting configuration is \textit{sensible}.
            \end{enumerate}
    \end{enumerate}
    \item \textbf{A \textit{stuck} configuration can never be reached from a \textit{sensible} configuration}. In fact we've proven that a transition from a \textit{sensible} configuration brings a another \textit{sensible} configuration (3).  We can prove by induction on the number of steps $n$ we can extend this property to every finite computation $\langle P,s\rangle \to^* \langle P_n, s_n\rangle$ (reflexive-transitive closure of $\to$)  to show that every reachable configuration starting from a \textit{sensible} configuration is \textit{sensible} too.
    \begin{enumerate}
        \item \textsc{Base case} ($n =0$)It's the same configuration.
        \item \textsc{Inductive case} ($n=k+1$) We know by Inductive Hypothesis that if $\langle P,s\rangle$ is \textit{sensible}, then $\langle P,s\rangle \to^*  \langle P_k, s_k\rangle$ brings to  $\langle P_k, s_k\rangle$ \textit{sensible} and by point 2 we know that it's not stuck. Then there is $\langle P_k, s_k\rangle \to \langle P_{k+1}, s_{k+1}\rangle $ which by point 3 it is \textit{sensible}.  
    \end{enumerate}
    We know that every \textit{sensible} configuration is not \textit{stuck} (2) hence every configuration reachable from a \textit{sensible} configuration is not \textit{stuck}.

\end{enumerate}
\end{solution}